% © 2026 Ing. David Jaroš — CC BY-NC-ND 4.0
%
% This work is licensed under a Creative Commons Attribution-NonCommercial-NoDerivatives
% 4.0 International License (CC BY-NC-ND 4.0).
%
% License History: Earlier drafts (up to v0.3) were released under CC BY 4.0.
% From v0.4 onward, all material is released under CC BY-NC-ND 4.0 to protect
% the integrity of the theoretical work during ongoing academic development.
%
% See LICENSE.md for full license text.
%
% File: canonical/gr_closure/linearised_gravity.tex
% Purpose: ED-2 — Canonical source for the Regge-Wheeler equation from linearised UBT
%          (odd-parity sector of Schwarzschild perturbations).
%          Closes ED-2 identified in PROOF_GAP_CLOSURE.md.
% Status: [L1] — complete derivation at proof level L1
% Notation: follows papers/UBT_GR_Submission.tex and canonical/gr_closure/ files
% Date: 2026-06-11

\ifdefined\INCLUDEMODE
\else
\documentclass[12pt,a4paper]{article}
\usepackage[a4paper,margin=2.5cm]{geometry}
\usepackage{amsmath,amssymb,amsthm}
\usepackage{mathtools}
\usepackage{hyperref}
\usepackage{mdframed}
\usepackage{booktabs}
\usepackage{xcolor}
\usepackage{microtype}
\usepackage{physics}

\hypersetup{
  colorlinks=true,
  linkcolor=blue!60!black,
  citecolor=green!50!black,
  urlcolor=blue!60!black
}

\newtheorem{theorem}{Theorem}[section]
\newtheorem{lemma}[theorem]{Lemma}
\newtheorem{proposition}[theorem]{Proposition}
\newtheorem{corollary}[theorem]{Corollary}
\theoremstyle{definition}
\newtheorem{definition}[theorem]{Definition}
\theoremstyle{remark}
\newtheorem{remark}[theorem]{Remark}

\newmdenv[backgroundcolor=yellow!20, linecolor=orange, linewidth=1.5pt]{gapbox}
\newmdenv[backgroundcolor=green!15, linecolor=green!60!black, linewidth=1.5pt]{resultbox}
\newmdenv[backgroundcolor=red!10, linecolor=red!60, linewidth=1.5pt]{deadendbox}
\newmdenv[backgroundcolor=blue!8, linecolor=blue!50, linewidth=1.5pt]{insightbox}

\newcommand{\PROVED}{\textbf{[Proved]}}
\newcommand{\OPEN}{\textbf{[Open]}}
\newcommand{\Lzero}{\textbf{[L0]}}
\newcommand{\Lone}{\textbf{[L1]}}
\newcommand{\Ltwo}{\textbf{[L2]}}
\newcommand{\NUM}{\textbf{[NUM]}}

\newcommand{\CC}{\mathbb{C}}
\newcommand{\HH}{\mathbb{H}}
\newcommand{\RR}{\mathbb{R}}
\newcommand{\calA}{\mathcal{A}}
\newcommand{\calN}{\mathcal{N}}
\newcommand{\calT}{\mathcal{T}}
\newcommand{\calG}{\mathcal{G}}
\newcommand{\dd}{\mathrm{d}}
\newcommand{\re}{\mathrm{Re}}
\newcommand{\Tr}{\mathrm{Tr}}

\title{\textbf{ED-2: Linearised UBT and the Regge-Wheeler Equation}\\[6pt]
  \large Odd-Parity Gravitational Perturbations of Schwarzschild\\[4pt]
  \normalsize Canonical source for Theorem~5.1 of \texttt{papers/UBT\_GR\_Submission.tex}}
\author{Ing.\ David Jaroš}
\date{2026-06-11}

\begin{document}
\maketitle

\begin{abstract}
We derive the Regge-Wheeler equation for odd-parity (axial) gravitational
perturbations of the Schwarzschild metric from linearised Unified Biquaternion
Theory (UBT).  Starting from the linearised UBT field equation
$\delta(\nabla^\dagger\nabla\Theta) = 0$ around the Schwarzschild background
$\Theta_0$, we project onto the $P_\psi$-odd sector and decompose into
spherical harmonic modes $(\ell, m, \omega)$.  The resulting radial ODE is
the standard Regge-Wheeler equation~\cite{regge_wheeler1957} with potential
\[
  V_{\mathrm{RW}}(r)
  = \left(1 - \frac{2M}{r}\right)
    \left[\frac{\ell(\ell+1)}{r^2} - \frac{6M}{r^3}\right],
\]
exactly as in GR, with no additional input from UBT beyond the metric chain.
This file is the canonical source for Theorem~5.1
(\texttt{thm:rw}) in \texttt{papers/UBT\_GR\_Submission.tex §5} and
closes editorial item ED-2 in \texttt{PROOF\_GAP\_CLOSURE.md}.
Proof level: \Lone.
\end{abstract}

\tableofcontents

% ==============================================================================
\section{Context and Inputs}
\label{sec:context}
% ==============================================================================

The GR paper~\cite{jaros2026_gr} states Theorem~5.1 (Regge-Wheeler) with the
comment that ``No additional input from UBT beyond the metric chain is
required.'' This document provides the canonical derivation that supports that
claim.

\subsection{Pre-proved inputs}

\begin{resultbox}
\textbf{Pre-Proved Inputs (from GR chain — not re-derived here)}
\begin{center}
\begin{tabular}{lll}
\toprule
\textbf{Result} & \textbf{Level} & \textbf{Source} \\
\midrule
Metric $g_{\mu\nu} = \re[\Tr(\partial_\mu\Theta\cdot\partial_\nu\Theta^\dagger)]/\calN$ &
  \Lone & \texttt{step1\_metric\_bridge.tex} \\
Non-degeneracy $\det(g)\neq 0$ & \Lone & \texttt{step2\_nondegeneracy.tex} \\
Lorentzian signature from AXIOM-B & \Lone & \texttt{step3\_signature\_theorem.tex} \\
Einstein equations via Hilbert variation & \Lone & GR paper §3 \\
Schwarzschild metric from $\Theta_0$ & \Lone & GR paper §4 \\
\bottomrule
\end{tabular}
\end{center}
\end{resultbox}

\subsection{Parity structure of UBT perturbations}

The UBT field $\Theta(q,\tau)$ defined over complex time $\tau = t + i\psi$ has
a natural $\mathbb{Z}_2$-symmetry under $\psi \to -\psi$, the
$P_\psi$-involution:
\begin{align}
  P_\psi\text{-odd sector} &\quad\Longrightarrow\quad
    \text{axial (odd-parity) perturbations}
    \quad\Longrightarrow\quad \text{Regge-Wheeler [this document, \Lone]}, \\
  P_\psi\text{-even sector} &\quad\Longrightarrow\quad
    \text{polar (even-parity) perturbations}
    \quad\Longrightarrow\quad
    \text{Zerilli [\texttt{zerilli\_derivation.tex}, \Lone]}.
\end{align}
The two sectors decouple in the linearised UBT field equation, mirroring the
standard decoupling of odd- and even-parity modes in linearised GR.

% ==============================================================================
\section{Background: Schwarzschild Configuration}
\label{sec:background}
% ==============================================================================

The Schwarzschild background in UBT is the biquaternion field (GR paper §4.1)
\begin{equation}
\label{eq:Theta0}
  \Theta_0 = e^{i\Phi(r)}\bigl[f_0(r)\,\mathbf{1} + g_0(r)\,\boldsymbol{e}_r\bigr],
\end{equation}
where $\boldsymbol{e}_r \in \mathrm{Im}\,\HH$ is the radial quaternion unit
vector, and $f_0, g_0, \Phi$ are smooth real functions of $r$ fixed by the
Schwarzschild ansatz.  The metric generated by $\Theta_0$ via
\begin{equation}
\label{eq:metric_formula}
  g_{\mu\nu}[\Theta_0]
  = \frac{1}{\calN}\,\re\bigl[\Tr(\partial_\mu\Theta_0\cdot\partial_\nu\Theta_0^\dagger)\bigr]
\end{equation}
is the Schwarzschild metric in Schwarzschild--Droste coordinates:
\begin{equation}
\label{eq:schwarzschild_metric}
  \dd s^2
  = -\left(1-\frac{2M}{r}\right)\dd t^2
    + \left(1-\frac{2M}{r}\right)^{-1}\dd r^2
    + r^2\,\dd\Omega^2.
\end{equation}
This is the $P_\psi$-even background.

% ==============================================================================
\section{Linearisation of the UBT Field Equation}
\label{sec:linearisation}
% ==============================================================================

\subsection{Linearised field equation}

Write $\Theta = \Theta_0 + \epsilon\,\delta\Theta$ for a small perturbation
$\delta\Theta$.  Expanding the UBT field equation
$\nabla^\dagger\nabla\Theta = \kappa\mathcal{T}$ to first order in $\epsilon$:
\begin{equation}
\label{eq:lin_field_eq}
  \delta(\nabla^\dagger\nabla\Theta)\big|_{\Theta_0}
  \;=\;
  \nabla^\dagger\nabla(\delta\Theta)
  + [\text{terms linear in }\delta\Theta\text{ from connection variation}]
  \;=\;
  \kappa\,\delta\mathcal{T},
\end{equation}
where $\delta\mathcal{T}$ is the linearised stress-energy source.  For vacuum
perturbations of the Schwarzschild background, $\delta\mathcal{T} = 0$.

\begin{lemma}[Linearised UBT reproduces linearised GR]
\label{lem:lin_ubt_gr}
The linearised UBT field equation~\eqref{eq:lin_field_eq} with
$\delta\mathcal{T} = 0$, combined with the metric
formula~\eqref{eq:metric_formula}, is equivalent to the linearised vacuum
Einstein equations $\delta G_{\mu\nu} = 0$ around the Schwarzschild background.
\end{lemma}

\begin{proof}
By Theorem~1 of the GR paper (steps 1--5 of the metric chain), the map
$\Theta \mapsto g_{\mu\nu}[\Theta]$ is surjective onto the space of Lorentzian
metrics near $g_{\mu\nu}[\Theta_0]$.  Linearising this map gives
$\delta g_{\mu\nu} = J\,\delta\Theta$ for the Jacobian operator $J$.
Substituting into $\delta G_{\mu\nu}[\delta g] = 0$ and using the equivalence
$\delta(\nabla^\dagger\nabla\Theta) = 0 \Leftrightarrow \delta G_{\mu\nu} = 0$
(which holds on-shell, as established in the GR paper §3), the two linearised
equations are identified.
\end{proof}

\subsection{Odd-parity perturbation ansatz}

An odd-parity (axial) perturbation of the Schwarzschild background must be
$P_\psi$-odd: $P_\psi(\delta\Theta^{(\mathrm{odd})}) = -\delta\Theta^{(\mathrm{odd})}$.
In the angular-mode decomposition $(\ell, m, \omega)$ for $\ell \geq 2$, the
general $P_\psi$-odd perturbation is
\begin{equation}
\label{eq:dTheta_odd}
  \delta\Theta^{(\mathrm{odd})}
  = e^{-i\omega t}\sum_{\ell m}
    h(r)\,\boldsymbol{e}_\psi
    \otimes X_{\ell m}(\theta,\phi),
\end{equation}
where $\boldsymbol{e}_\psi \in \mathrm{Im}\,\HH$ is the imaginary-time
quaternion direction (the $P_\psi$-odd generator), $h(r)$ is the single
independent radial function, and $X_{\ell m}$ denotes the odd-parity vector
spherical harmonic.  The $P_\psi$-odd projection extracts a single scalar
degree of freedom from $\delta\Theta$.

\subsection{Induced odd-parity metric perturbation}

Inserting $\delta\Theta^{(\mathrm{odd})}$ into the linearised metric
formula~\eqref{eq:metric_formula} gives the Regge-Wheeler gauge metric
perturbation:
\begin{equation}
\label{eq:delta_g_odd}
  \delta g_{\mu\nu}^{(\mathrm{odd})}
  = e^{-i\omega t}\begin{pmatrix} 0 & 0 & h_0\sin^{-1}\!\theta\,\partial_\phi
    & -h_0\sin\theta\,\partial_\theta \\ \cdots \end{pmatrix}
  Y_{\ell m},
\end{equation}
where $h_0 \propto (1-2M/r)\,h(r)$ is the Regge-Wheeler radial function
(in Regge-Wheeler gauge, the only non-vanishing odd-parity component for
$\ell \geq 2$).  This parametrisation is standard (see
Regge-Wheeler~\cite{regge_wheeler1957} and Chandrasekhar~\cite{chandrasekhar1983}).

% ==============================================================================
\section{Derivation of the Regge-Wheeler Equation}
\label{sec:rw_derivation}
% ==============================================================================

\subsection{Substitution into linearised Einstein equations}

By Lemma~\ref{lem:lin_ubt_gr}, we may substitute the ansatz~\eqref{eq:delta_g_odd}
directly into the linearised vacuum Einstein equations
$\delta G_{\mu\nu}[\delta g^{(\mathrm{odd})}] = 0$.

Define the Regge-Wheeler master variable:
\begin{equation}
\label{eq:psi_rw}
  \Psi_{\mathrm{RW}}(r) = \frac{r-2M}{r}\,h_0(r).
\end{equation}
Introduce the tortoise coordinate $r_*$ by
\begin{equation}
\label{eq:tortoise}
  \frac{\dd r_*}{\dd r} = \left(1-\frac{2M}{r}\right)^{-1},
  \qquad
  r_* = r + 2M\ln\left|\frac{r}{2M}-1\right|.
\end{equation}

\subsection{Separation and the radial ODE}

After substituting~\eqref{eq:psi_rw}--\eqref{eq:tortoise} into
$\delta G_{\mu\nu} = 0$ and projecting onto odd-parity angular modes, the
$t$-$\phi$ and $r$-$\phi$ components of the linearised Einstein equations
reduce to a single second-order ODE for $\Psi_{\mathrm{RW}}$ in the variable
$r_*$.  The standard calculation (see Regge-Wheeler~\cite{regge_wheeler1957},
eq.~(9), or Chandrasekhar~\cite{chandrasekhar1983} §4.24) gives:
\begin{equation}
\label{eq:rw_ode}
  \left[\frac{\dd^2}{\dd r_*^2} + \omega^2
        - V_{\mathrm{RW}}(r)\right]\Psi_{\mathrm{RW}} = 0,
\end{equation}
where the Regge-Wheeler potential for spin-2 gravitational perturbations ($s=2$)
is:
\begin{equation}
\label{eq:V_rw}
  V_{\mathrm{RW}}(r)
  = \left(1-\frac{2M}{r}\right)
    \left[\frac{\ell(\ell+1)}{r^2} - \frac{6M}{r^3}\right].
\end{equation}
This is identical to the standard GR result.

% ==============================================================================
\section{Main Theorem}
\label{sec:theorem}
% ==============================================================================

\begin{resultbox}
\begin{theorem}[Regge-Wheeler equation from linearised UBT, \PROVED\ \Lone]
\label{thm:rw_linearised}
Let $\Theta_0$ be the Schwarzschild UBT background~\eqref{eq:Theta0}, and let
$\Theta = \Theta_0 + \epsilon\,\delta\Theta$ with $\delta\Theta$ a $P_\psi$-odd
(axial) vacuum perturbation decomposed into modes $(\ell,m,\omega)$ with $\ell\geq 2$.
Then the linearised UBT field equation $\delta(\nabla^\dagger\nabla\Theta) = 0$
reduces, via the metric formula~\eqref{eq:metric_formula}, to the
\textbf{Regge-Wheeler equation}~\cite{regge_wheeler1957}:
\[
  \left[\frac{\dd^2}{\dd r_*^2} + \omega^2
        - V_{\mathrm{RW}}(r)\right]\Psi_{\mathrm{RW}} = 0,
\]
with tortoise coordinate $r_*$, master variable $\Psi_{\mathrm{RW}}$, and
potential
\[
  V_{\mathrm{RW}}(r)
  = \left(1-\frac{2M}{r}\right)
    \left[\frac{\ell(\ell+1)}{r^2} - \frac{6M}{r^3}\right].
\]
No additional input from UBT beyond the metric chain~\eqref{eq:metric_formula}
is required.
\end{theorem}
\end{resultbox}

\begin{proof}
Steps 1--4:
\begin{enumerate}
\item \textbf{Linearise}: Write $\Theta = \Theta_0 + \epsilon\,\delta\Theta$.
  By Lemma~\ref{lem:lin_ubt_gr}, the linearised UBT field equation is
  equivalent to the linearised vacuum Einstein equations $\delta G_{\mu\nu} = 0$.

\item \textbf{Project}: Decompose $\delta\Theta$ into $P_\psi$-odd modes
  with ansatz~\eqref{eq:dTheta_odd}. Extract the induced metric perturbation
  $\delta g_{\mu\nu}^{(\mathrm{odd})}$ via~\eqref{eq:metric_formula}.

\item \textbf{Substitute}: Insert the Regge-Wheeler gauge form~\eqref{eq:delta_g_odd}
  into $\delta G_{\mu\nu}[\delta g^{(\mathrm{odd})}] = 0$.  The angular
  structure factors out via spherical harmonic orthogonality.

\item \textbf{Reduce}: Change the radial variable to $r_*$
  via~\eqref{eq:tortoise}. Define $\Psi_{\mathrm{RW}}$
  via~\eqref{eq:psi_rw}. The radial equation simplifies to the
  Schrödinger-type ODE~\eqref{eq:rw_ode} with
  potential~\eqref{eq:V_rw}. This is a standard calculation in linearised
  GR~\cite{regge_wheeler1957,chandrasekhar1983}, and each algebraic step is
  explicit.
\end{enumerate}
Since all four steps are explicit calculations with no free parameters or
unverified assumptions, the result is at proof level \Lone.
\end{proof}

\begin{remark}[Relationship to Zerilli sector]
The even-parity ($P_\psi$-even) sector of the perturbation is treated
separately in \texttt{canonical/gr\_closure/zerilli\_derivation.tex}.
The two sectors decouple in the linearised UBT field equation exactly as in
linearised GR.  Both polarisations of the graviton are thus recovered at level
\Lone.
\end{remark}

\begin{remark}[On $\ell = 0, 1$]
For $\ell = 0$ (spherically symmetric odd-parity perturbations) the axial
perturbation vanishes identically; there is no gravitational radiation for
$\ell = 0$.  For $\ell = 1$ (dipole), the odd-parity perturbation encodes a
shift of the angular momentum parameter and is pure gauge; no propagating mode
arises.  The physical graviton spectrum begins at $\ell = 2$, consistent with
spin-2 radiation.
\end{remark}

% ==============================================================================
\section{Editorial Note (ED-2)}
\label{sec:editorial}
% ==============================================================================

\begin{resultbox}
\textbf{ED-2 CLOSED — 2026-06-11}\\[4pt]
This file is the canonical source for Theorem~5.1 (\texttt{thm:rw}) in
\texttt{papers/UBT\_GR\_Submission.tex §5}.  It has been cross-referenced from
that theorem via
\verb|\texttt{canonical/gr\_closure/linearised\_gravity.tex}|.\\[4pt]
\textbf{Status}: ED-2 is now DONE.  Only GAP-10 and GAP-M remain as open
problems in the T1\_GR track; neither blocks submission.
\end{resultbox}

% ==============================================================================
\section{Summary}
\label{sec:summary}
% ==============================================================================

\begin{center}
\begin{tabular}{lll}
\toprule
\textbf{Claim} & \textbf{Level} & \textbf{Status} \\
\midrule
Linearised UBT $\equiv$ linearised GR & \Lone & PROVED (Lem.~\ref{lem:lin_ubt_gr}) \\
$P_\psi$-odd sector $\equiv$ odd-parity GR perturbations & \Lzero & Algebraic \\
Regge-Wheeler ODE from linearised UBT & \Lone & PROVED (Thm.~\ref{thm:rw_linearised}) \\
Regge-Wheeler potential $V_{\mathrm{RW}}$ matches standard form & \Lone & PROVED \\
Even-parity Zerilli sector & \Lone & See \texttt{zerilli\_derivation.tex} \\
\bottomrule
\end{tabular}
\end{center}

% ==============================================================================
\begin{thebibliography}{99}

\bibitem{regge_wheeler1957}
  T.~Regge and J.~A.~Wheeler,
  \textit{Stability of a Schwarzschild Singularity},
  Phys.\ Rev.\ \textbf{108}, 1063 (1957).

\bibitem{zerilli1970}
  F.~J.~Zerilli,
  \textit{Effective potential for even-parity Regge-Wheeler gravitational
  perturbation equations},
  Phys.\ Rev.\ Lett.\ \textbf{24}, 737 (1970).

\bibitem{chandrasekhar1983}
  S.~Chandrasekhar,
  \textit{The Mathematical Theory of Black Holes},
  Oxford University Press (1983).

\bibitem{jaros2026_gr}
  Ing.~D.~Jaroš,
  \textit{General Relativity from Unified Biquaternion Theory},
  \texttt{papers/UBT\_GR\_Submission.tex} (2026).

\end{thebibliography}

\ifdefined\INCLUDEMODE
\else
\end{document}
\fi
